\section{课后练习}

\subsection{代码质量}

\subsubsection{在 pyproject.toml 中配置 ruff 规则}
ruff 的规则与格式化参数可在 pyproject.toml 中集中配置，避免在命令行反复指定。
\begin{minted}[frame=single, fontsize=\small]{ini}
[tool.ruff]
line-length = 88
target-version = "py313"
[tool.ruff.lint]
select = ["E", "F", "I", "W"]
[tool.pytest.ini_options]
testpaths = ["."]
python_files = "test_*.py"
\end{minted}

\subsubsection{用 pytest 参数化测试减少重复}
使用 pytest.mark.parametrize 可以对同一测试函数批量提供多组输入，避免复制粘贴。
\begin{minted}[frame=single, fontsize=\small]{python}
import pytest

@pytest.mark.parametrize("name", ["Alice", "Bob", "Carol"])
def test_normal_names(name, capsys, monkeypatch):
    monkeypatch.setattr("sys.argv", ["sdt-greet", "--name", name])
    main()
    assert f"Hello, {name}!" in capsys.readouterr().out
\end{minted}

\subsubsection{用 pytest.raises 断言异常与退出码}
对会抛出异常的代码，用 pytest.raises 既能捕获异常，又能断言退出码，是测试 CLI 参数校验的常用方式。
\begin{minted}[frame=single, fontsize=\small]{python}
import pytest

def test_blank_name(monkeypatch):
    monkeypatch.setattr(sys, "argv", ["sdt-greet", "--name", "   "])
    with pytest.raises(SystemExit) as exc_info:
        main()
    assert exc_info.value.code == 2
\end{minted}

\subsection{元编程：构建系统}

\subsubsection{GNU Make 的自动变量}
自动变量 \verb|$@|、\verb|$<|、\verb|$^| 分别代表目标文件、第一个依赖和全部依赖，可让配方更通用。
\begin{minted}[frame=single, fontsize=\small]{make}
output.txt: input.txt
	cat $< > $@      # $< 即 input.txt，$@ 即 output.txt
\end{minted}

\subsubsection{用 .PHONY 声明伪目标}
当目标名与文件名同名时，make 可能误判其已是最新；用 .PHONY 声明后，每次都会执行该配方。
\begin{minted}[frame=single, fontsize=\small]{make}
.PHONY: all clean build
all: build
clean:
	rm -rf dist/ *.egg-info
\end{minted}

\subsection{大杂烩：API、jq 与 Markdown}

\subsubsection{curl 的常用选项}
-f 在 HTTP 错误时静默失败，-s 关闭进度条，-o 保存到文件，三者组合适合脚本化的 API 调用。
\begin{minted}[frame=single, fontsize=\small]{bash}
curl -fsS http://127.0.0.1:8000/packages.json
curl -fsS http://127.0.0.1:8000/packages.json -o packages.json
\end{minted}

\subsubsection{用 jq 筛选与排序 JSON}
jq 的 select 用于按条件过滤，sort\_by 可按一个或多个键排序，负号表示降序。
\begin{minted}[frame=single, fontsize=\small]{bash}
jq '[.[] | select(.status == "active" and .downloads >= 100)]
    | sort_by([-.downloads, .name])'
\end{minted}

\subsubsection{用 jq -r 拼接 Markdown 表格}
-j 使用原始输出，配合字符串拼接可直接生成 Markdown 表格行。
\begin{minted}[frame=single, fontsize=\small]{bash}
jq -r '.[] | "| \(.name) | \(.version) | \(.downloads) |"'
\end{minted}

\subsection{Python 与 CLI 约定}

\subsubsection{用 argparse 定义必选参数}
argparse 的 required=True 可声明必选参数，--help 由框架自动生成，符合常见 CLI 约定。
\begin{minted}[frame=single, fontsize=\small]{python}
import argparse

p = argparse.ArgumentParser()
p.add_argument("--name", required=True)
a = p.parse_args()
print(f"Hello, {a.name}!")
\end{minted}

\subsubsection{计算文件的 SHA-256 校验和}
使用 sha256sum 或 Python 的 hashlib 计算文件摘要，可用于校验构建产物的完整性。
\begin{minted}[frame=single, fontsize=\small]{bash}
sha256sum dist/greetlab_25020007136-0.1.0-py3-none-any.whl
\end{minted}

\begin{minted}[frame=single, fontsize=\small]{python}
import hashlib
h = hashlib.sha256(open("file.whl", "rb").read())
print(h.hexdigest())
\end{minted}
